\documentclass[exercice,a4]{thedocument}%
\usepackage[french]{babel}%
\usepackage{tabularx}%
\usepackage{array}%
\startdatakeys
    \source{21GENMATNC1}
    \BO{2026}
    \cycle{4}
    \competencesDuSocle{CA,RA}
    \objectifsApprentissage{C4C1bc,C4C2bd}
    \motsCles{probabilité}
\enddatakeys
\begin{document}%

\medskip

Un jeu est constitué de quatre familles de cartes:
banane ;
prune ;
citron ;
fraise.

Voici la répartition des cartes de la famille banane.

\begin{center}
\begin{tabularx}{\linewidth}{|m{3.5cm}|*{5}{>{\centering \arraybackslash}X|}}\hline
Nombre de banane(s)	&1&2&3&4&5\\ \hline
Nombre de cartes	&5&3&3&2&1\\ \hline
\end{tabularx}
\end{center}

La répartition est la même pour les cartes avec les autres fruits.

\medskip

\begin{enumerate}
\item Montrer que ce jeu a $56$ cartes.

Joanne mélange toutes les cartes. Son frère Jack prend une carte au hasard. On admet que chaque carte a la même chance d'être choisie.

Soit $P$ l'évènement: \og Jack obtient une carte de la famille prune \fg.
\item Quelle est la probabilité de l'évènement $P$ ?
\item
	\begin{enumerate}
		\item Quel est l'évènement contraire de $P$ ?
		\item Quelle est la probabilité de l'évènement contraire de $P$ ?
	\end{enumerate}
\item Quelle est la probabilité d'obtenir une carte avec quatre fruits ?
\end{enumerate}

\bigskip

\showthesource
\end{document}
